\documentclass[10pt,a4paper]{article}
\usepackage[T1]{fontenc}
\usepackage{lmodern}
\usepackage[left=17mm,right=17mm,top=15mm,bottom=16mm]{geometry}
\usepackage{amsmath,booktabs,tabularx,array}
\usepackage[hidelinks]{hyperref}
\setlength{\parindent}{0pt}
\setlength{\parskip}{3pt}
\setlength{\tabcolsep}{4pt}
\renewcommand{\arraystretch}{1.12}
\newcolumntype{Y}{>{\raggedright\arraybackslash}X}
\newcommand{\sect}[1]{\vspace{5pt}\textbf{#1}\par\vspace{1pt}}
\newcommand{\model}[1]{\texttt{#1}}
\pagestyle{plain}
\begin{document}
\begin{center}
{\small Special Topics: Social Media Mining and Analysis -- Assignment 1}\par
\vspace{3pt}
{\Large\bfseries Unimodal Baselines for Hateful Meme Classification}\par
\vspace{5pt}
Bogdan Fanea\qquad Lena Bryll\qquad Natalie Stadler\par
\vspace{2pt}
{\small\url{https://github.com/BogdanFanea/SpecialTopicSocialMining}}
\end{center}
\vspace{-3pt}

We tested four pretrained encoders on the Hateful Memes dataset~\cite{kiela}, using
either images or the supplied meme text. Each encoder stays frozen while a separate
linear classifier learns from its features. All four classifiers use the same design.

\sect{1. Data and evaluation labels}
\begin{tabularx}{\linewidth}{@{}Yrrr@{}}
\toprule
Split & Examples & Non-hateful (0) & Hateful (1)\\
\midrule
Training & 8,500 & 5,450 & 3,050\\
Validation (\model{dev}) & 500 & 250 & 250\\
Test (public-mirror labels) & 1,000 & 510 & 490\\
\bottomrule
\end{tabularx}

The splits have no shared IDs, and all images are present. We kept the supplied
training and validation labels. The course test file has no labels, so we used a pinned
public mirror~\cite{labels} with the same 1,000 IDs, texts, and image paths. Its test
annotation file is byte-identical to the Kaggle copy cited by a research implementation.
\textbf{Official ground-truth provenance is unconfirmed.} Label versions differ:
the mirror changes 823 training and three validation labels; we apply neither
change. Test scores below are therefore \emph{AUROC against public-mirror labels},
not authenticated competition scores. Source revision and checksums are recorded
in the accompanying provenance file.

\sect{2. Encoders and preprocessing}
\begin{tabularx}{\linewidth}{@{}lYr l@{}}
\toprule
Baseline & Exact pretrained checkpoint & Dim. & Pooling\\
\midrule
Image A & \model{ResNet50\_Weights.IMAGENET1K\_V2} & 2,048 & Global average\\
Image B & \model{openai/clip-vit-base-patch32} & 512 & Vision CLS + projection\\
Text A & \model{sentence-transformers/all-MiniLM-L6-v2} & 384 & Masked mean + L2\\
Text B & \model{sentence-transformers/all-mpnet-base-v2} & 768 & Masked mean + L2\\
\bottomrule
\end{tabularx}

We chose ResNet-50 as an ImageNet-trained CNN baseline and CLIP as a vision
transformer trained with image--text pairs~\cite{vision,clip}. MiniLM is a compact
six-layer sentence encoder; MPNet is a larger twelve-layer model~\cite{textmodels}.
Because these encoders differ in architecture, size, and pretraining, the results
cannot isolate the effect of any one of these factors.

Image branches receive \textbf{RGB pixels only}. ResNet uses its official bilinear
shorter-edge resize to 232 and $224\times224$ center crop with ImageNet channel
normalization. CLIP uses the checkpoint's bicubic resize and $224\times224$ center
crop with CLIP channel normalization; its text tower is unused. Text branches
receive \textbf{only the provided meme text}, using their matching tokenizers,
dynamic padding, and truncation limits of 256 (MiniLM) and 384 (MPNet) tokens,
including special tokens. Attention-mask mean pooling excludes padding, includes
non-padding special tokens, and is followed by L2 normalization. No text was
truncated in any split. Encoders run in evaluation/inference mode; float32
embeddings are cached once. Exact model revisions and normalization constants
are documented in \model{config.json} and the README.

\sect{3. Linear probes and model selection}
For every encoder we fit $p(y=1\mid z)=\sigma(w^\top\widetilde z+b)$, where
$\widetilde z$ uses per-dimension mean/scale fitted \textbf{only on training
embeddings}. The classifier is \model{Linear(d,1)} with bias and mean
\model{BCEWithLogitsLoss}; only $d$ varies. AdamW uses batch size 256, 50 epochs,
learning rates $\{10^{-4},10^{-3},10^{-2}\}$ and weight decay $\{0,0.01\}$
(including bias), with default betas/epsilon and no scheduler or class weighting.
We ran all six configurations with seeds $\{42,43,44\}$: 72 fits in total. For each
fit, we kept the epoch with the highest validation AUROC, then chose settings by mean
validation AUROC across seeds. Ties keep the earliest epoch or first grid entry.
We evaluated on test after selection; validation and test data were never used for training.

Selected settings (learning rate, weight decay; best epochs for seeds 42/43/44):
ResNet $(.001,.01;\ 1/2/1)$; CLIP $(.01,.01;\ 4/10/9)$;
MiniLM $(.001,0;\ 5/8/10)$; MPNet $(.001,.01;\ 5/2/1)$.
Extraction used Apple M3 MPS (16 GB); deterministic, single-thread CPU probe
training used PyTorch 2.4.0 and Transformers 4.44.2.

\begin{center}
\footnotesize
\begin{tabular}{@{}c c c c c c c c c@{}}
\fbox{Pixels} & $\to$ & \fbox{Frozen ResNet / CLIP} & $\to$ & \fbox{Cached $z$} & $\to$ & \fbox{Train-scaled linear probe} & $\to$ & \fbox{Sigmoid}\\[4pt]
\fbox{Meme text} & $\to$ & \fbox{Frozen MiniLM / MPNet} & $\to$ & \fbox{Cached $z$} & $\to$ & \fbox{Train-scaled linear probe} & $\to$ & \fbox{Sigmoid}
\end{tabular}\par
\vspace{2pt}
Figure 1. Four independent branches; the alternatives in each row are run separately.
\end{center}

\newpage
\sect{4. Results and comparison}
Table 1 reports AUROC on the 0--1 scale, computed from continuous sigmoid
hateful-class scores. Values are mean $\pm$ sample standard deviation over the
same three seeds; predictions are not ensembled.

\begin{tabularx}{\linewidth}{@{}lYcc@{}}
\toprule
Baseline & Encoder checkpoint & Validation AUROC & Test AUROC$^\dagger$\\
\midrule
Image A & \model{ResNet50\_Weights.IMAGENET1K\_V2} & $0.5429\pm0.0218$ & $0.5390\pm0.0091$\\
Image B & \model{openai/clip-vit-base-patch32} & $\mathbf{0.6445}\pm0.0049$ & $0.6552\pm0.0033$\\
Text A & \model{sentence-transformers/all-MiniLM-L6-v2} & $0.6310\pm0.0019$ & $\mathbf{0.6603}\pm0.0016$\\
Text B & \model{sentence-transformers/all-mpnet-base-v2} & $0.6390\pm0.0031$ & $0.6566\pm0.0069$\\
\bottomrule
\end{tabularx}
{\footnotesize\textbf{Table 1.} $^\dagger$AUROC against public-mirror labels; see Section 1.
Bold marks the highest mean in each column.}

CLIP scored higher than ResNet-50 on both splits. MPNet was the stronger text
model on validation, but MiniLM had the higher test mean, so the larger text encoder
did not perform better on both splits. We selected CLIP and MPNet for the case
analysis using validation scores.
Their test means, 0.6552 and 0.6566, differ by only 0.0014. Even MiniLM's highest
observed test mean exceeds CLIP by just 0.0051. These small differences do not
establish a reliable modality advantage: three seeds measure optimizer variability,
not sampling confidence, and validation estimates were reused for selection.

CLIP's training on image--text pairs may help it capture meaning beyond ImageNet
object categories, though this comparison cannot establish the cause. Text models
can use words associated with hate while missing what an image refers to or whether
a caption rejects a hateful message. Confounder pairs challenge both input types.
\textbf{Image models still receive the written captions as pixels}, and CLIP's
pretraining used text. Center cropping may also remove useful context at image edges.

\sect{5. Five validation disagreements}
We compare CLIP and MPNet using seed 42 and the fixed rule
\textbf{hateful iff score $\geq0.5$}. They disagree on 167 of 500 validation
examples. We inspect the five disagreements with the largest absolute score
gap (ties by ascending ID). In all five cases, CLIP predicts non-hateful and MPNet predicts hateful.
These selected examples show possible errors, but not how common each error is.
AUROC itself uses no threshold. Labels below are the supplied course labels.

\begin{tabularx}{\linewidth}{@{}rcccY@{}}
\toprule
ID & True $y$ & $p_\mathrm{image}$ & $p_\mathrm{text}$ & Inspected context and likely error\\
\midrule
38047 & 0 & 0.160 & 0.853 & A gendered insult overlays a political group photo.
MPNet flags offensive wording; the supplied non-hateful label is ambiguous,
highlighting the distinction between offensiveness and group targeting.\\[3pt]
82509 & 0 & 0.145 & 0.818 & A runway photo accompanies a vulgar ``72 virgins''
caption. MPNet may overreact to provocative wording without resolving the target;
the supplied label is non-hateful.\\[3pt]
01967 & 0 & 0.118 & 0.765 & A burning swastika resolves the caption's reference
to racism's ``symbol.'' MPNet's false positive may miss the anti-racist stance
and the visual referent.\\[3pt]
48296 & 1 & 0.156 & 0.771 & A woman playfully bites a dog's ear; a caption about
an Asian boyfriend's family invokes a racist dog-eating stereotype.
CLIP misses this image--text connection.\\[3pt]
19536 & 1 & 0.052 & 0.664 & An altered-nose portrait and a caption about a
Jewish friend finding coins invoke antisemitic stereotypes. CLIP misses the
combination; MPNet can exploit explicit lexical cues.\\
\bottomrule
\end{tabularx}

These explanations come from inspecting the images and captions; we did not test
which features drove each decision. A correct text prediction does not show that a
model understood a joke that depends on the image. Ambiguous labels and the uncertain
source of the test labels also limit the conclusions. We did not change any labels.

\sect{Sources and reused implementations}
The code uses TorchVision/Hugging Face model loaders, PyTorch optimization, and
scikit-learn scaling/metrics. Masked pooling follows the text model-card examples.
Configurations, all selected
classifier checkpoints, per-seed predictions, and complete search histories accompany
the submission.
{\small
\begin{thebibliography}{5}
\setlength{\itemsep}{0pt}\setlength{\parskip}{0pt}
\bibitem{kiela} D. Kiela et al. \href{https://arxiv.org/abs/2005.04790}{The Hateful Memes Challenge: Detecting Hate Speech in Multimodal Memes}. NeurIPS, 2020.
\bibitem{vision} TorchVision. \href{https://docs.pytorch.org/vision/0.19/models/generated/torchvision.models.resnet50.html}{ResNet-50 pretrained weights and inference transforms}, version 0.19.
\bibitem{clip} A. Radford et al. \href{https://arxiv.org/abs/2103.00020}{Learning Transferable Visual Models From Natural Language Supervision}. ICML, 2021. \href{https://huggingface.co/openai/clip-vit-base-patch32}{CLIP checkpoint}.
\bibitem{textmodels} Sentence Transformers. Model cards and pooling examples for \href{https://huggingface.co/sentence-transformers/all-MiniLM-L6-v2}{all-MiniLM-L6-v2} and \href{https://huggingface.co/sentence-transformers/all-mpnet-base-v2}{all-mpnet-base-v2}.
\bibitem{labels} \href{https://huggingface.co/datasets/neuralcatcher/hateful_memes/blob/d201c488dc7024623d1ecbcc987b3f132c4c2e12/test_seen.jsonl}{Public test-label mirror, revision d201c488}; \href{https://github.com/YiLunLee/missing_aware_prompts/blob/9aa538159f4623de1936d25daee880ec47c062be/DATA.md}{Missing-Aware Prompts data instructions}. Accessed 17 October 2026.
\end{thebibliography}
}
\end{document}
